\documentclass[12pt]{article}
\usepackage{amsmath, amssymb}
\usepackage{geometry}
\usepackage{hyperref}
\usepackage{parskip}
\usepackage{microtype}
\usepackage{booktabs}
\usepackage{xcolor}
\geometry{margin=1.3in}

\definecolor{gold}{RGB}{180,140,30}

\title{\textbf{Generation Mixing and CP Violation\\
from Channel Weights}}
\author{Joseph Shields}
\date{2026}

\begin{document}
\maketitle

\begin{abstract}
The single recursion $c^2 = c+1$ fixes the contraction rate $r = 1/(2\varphi)$
and the three channel weights $W_L$, $W_B$, $W_M$. From these, three
fundamental CKM observables are derived with no free parameters.
The Wolfenstein mixing parameter follows from the algebraic identity
$r/(1+r) = 1/\varphi^3$, giving $\lambda = 0.23607$ (observed $|V_{us}|
= 0.22450$, residual $+5.15\% = 0.58\,n\,\varepsilon_{\text{floor}}$, $n=3$).
The CP-violating phase follows from the angle whose cosine equals the
boundary-channel weight: $\delta_{\text{CP}} = \arccos(W_B) = 1.1296$\,rad
(PDG 2022: $1.14 \pm 0.13$\,rad, residual $-0.92\% = 0.31\,\varepsilon_{\text{floor}}$).
The $b \to c$ transition amplitude $|V_{cb}| = W_B \times W_M = 0.04078$
(PDG $0.04113$, residual $-0.85\% = 0.14\,n\,\varepsilon_{\text{floor}}$, $n=2$)
derives the Wolfenstein parameter $A = 0.7318$ (PDG $0.7380$) from the
same channel weights, completing the leading CKM sector.
\end{abstract}

%% ─────────────────────────────────────────────────────────────────────────────
\section{Foundations}
%% ─────────────────────────────────────────────────────────────────────────────

The axiom and its consequences:
\begin{equation}
  \varphi^2 = \varphi + 1
  \quad\Rightarrow\quad
  \varphi = \tfrac{1+\sqrt{5}}{2},\qquad
  r = \tfrac{1}{2\varphi} \approx 0.30902.
  \label{eq:axiom}
\end{equation}

The channel weights and braiding floor:
\begin{align}
  W_L &= (1-r)^2     \approx 0.4775, \label{eq:wl}\\
  W_B &= 2r(1-r)     \approx 0.4271, \label{eq:wb}\\
  W_M &= r^2         \approx 0.0955, \label{eq:wm}\\
  \varepsilon_{\text{floor}} &= r^3 \approx 0.02951. \label{eq:eps}
\end{align}

In the UM picture, matter propagates through three channels: a
longitudinal channel (dominant, weight $W_L$) carrying the bulk of
momentum flow; a boundary channel (weight $W_B$) at the interface between
successive recursion cells, carrying both electromagnetic and weak
interactions; and a matter channel (weight $W_M$) associated with Higgs
coupling and mass generation. Generation structure arises when quark-sector
trajectories braid across cell boundaries, and inter-generation transitions
occur precisely at those boundary crossings.

%% ─────────────────────────────────────────────────────────────────────────────
\section{The Wolfenstein Parameter}
%% ─────────────────────────────────────────────────────────────────────────────

\subsection{Physical origin}

The CKM matrix parametrises transitions between quark generations. In the
Standard Model its entries are free parameters fitted to experiment. In UM,
quark generations correspond to distinct recursion-cell levels, and
inter-generation transitions are boundary-crossing events. A quark of
generation $n$ transitions to generation $n \pm 1$ by traversing one
boundary crossing; the amplitude for this transition is proportional to the
boundary-to-total ratio of the contraction rate.

The relevant amplitude is
\begin{equation}
  \lambda = \frac{r}{1 + r},
  \label{eq:lambda_raw}
\end{equation}
where $r$ is the contraction rate and $1 + r$ is the total (interior +
boundary) channel factor.

\subsection{Algebraic identity: $r/(1+r) = 1/\varphi^3$}

This expression simplifies exactly using the defining property
$\varphi^2 = \varphi + 1$:

\begin{align}
  \frac{r}{1+r}
  &= \frac{1/(2\varphi)}{1 + 1/(2\varphi)}
   = \frac{1}{2\varphi + 1}. \label{eq:id1}
\end{align}

Now $2\varphi + 1 = \varphi + \varphi + 1 = \varphi + \varphi^2 =
\varphi(1 + \varphi) = \varphi \cdot \varphi^2 = \varphi^3$, using
$\varphi^2 = \varphi + 1$ twice. Therefore
\begin{equation}
  \boxed{\frac{r}{1+r} = \frac{1}{\varphi^3}}.
  \label{eq:id}
\end{equation}

This identity is exact; it holds for any $\varphi$ satisfying $\varphi^2 =
\varphi + 1$. It is not a numerical coincidence: it is a consequence of the
axiom $\varphi^2 = \varphi + 1$ applied to the ratio of channel amplitudes.

\subsection{Numerical result}

\begin{equation}
  \lambda_{\text{UM}} = \frac{1}{\varphi^3}
  = \frac{1}{(1.61803\ldots)^3}
  = 0.23607.
  \label{eq:lambda_num}
\end{equation}

The observed value is $|V_{us}| = 0.22450$ (PDG 2022). The residual is
\begin{equation}
  \text{Residual} = \frac{0.23607 - 0.22450}{0.22450}
  = +5.15\%.
  \label{eq:lambda_res}
\end{equation}

The expression $\lambda = 1/\varphi^3$ is a 3-cycle quantity; its accumulated
braiding floor is $n\,\varepsilon_{\text{floor}} = 3 \times r^3 \approx 8.85\%$.
Normalised by the accumulated floor:
\[
  \frac{5.15\%}{8.85\%} = 0.58\,n\,\varepsilon_{\text{floor}}.
\]
The $\lambda$ result is therefore within the accumulated floor. The
QCD running of quark-sector trajectories is expected to supply the residual
correction; its derivation is reserved for the next paper in this series.

%% ─────────────────────────────────────────────────────────────────────────────
\section{The CP-Violating Phase}
%% ─────────────────────────────────────────────────────────────────────────────

\subsection{Physical origin}

CP violation requires a complex relative phase between interfering
amplitudes. In UM, the boundary channel $W_B$ carries both the
electromagnetic and weak interactions; it is the locus of inter-generation
transitions. When two such transitions interfere, the relative phase of
their amplitudes on the unit circle is the physical CP-violating angle.

The boundary channel has weight $W_B = 2r(1-r)$. This weight is the
projection of the boundary-channel amplitude onto the real axis of the
unit circle in the channel-weight space. The angle subtended by this
projection is
\begin{equation}
  \delta_{\text{CP}} = \arccos(W_B) = \arccos\!\bigl(2r(1-r)\bigr).
  \label{eq:delta}
\end{equation}

No additional parameter is introduced. The CP phase is the geometric angle
in channel-weight space whose cosine equals $W_B$.

\subsection{Numerical result}

\begin{equation}
  W_B = 2 \times 0.30902 \times 0.69098 = 0.42705,
\end{equation}
\begin{equation}
  \delta_{\text{CP}}
  = \arccos(0.42705)
  = 1.1296\,\text{rad}
  = 64.72^\circ.
  \label{eq:delta_num}
\end{equation}

The PDG 2022 world-average value from global CKM fits is
$\delta_{\text{CP}} = 1.14 \pm 0.13$\,rad. The residual is

\begin{equation}
  \text{Residual}
  = \frac{1.1296 - 1.14}{1.14}
  = -0.92\% = -0.31\,\varepsilon_{\text{floor}}.
  \label{eq:delta_res}
\end{equation}

\textcolor{gold}{This is the cleanest new result of the paper.} At
$0.31\,\varepsilon_{\text{floor}}$, the UM-derived CP phase lies well
within the sub-floor regime; no correction cycle is geometrically mandated.
It is also the first time a CP-violating phase has been derived from a
single geometric axiom rather than measured and inserted as a free
parameter.

\subsection{Comparison with the PDG uncertainty}

The PDG 1$\sigma$ uncertainty on $\delta_{\text{CP}}$ is $\pm 0.13$\,rad,
or $\pm 11\%$. The UM residual of $0.92\%$ is twelve times smaller than
the experimental uncertainty. Within the current precision of global CKM
fits, the UM value $\arccos(W_B)$ is indistinguishable from the measured
value.

As experimental precision on $\delta_{\text{CP}}$ improves from CP-violation
experiments (Belle~II, LHCb), this will become one of the sharpest tests
of UM: the framework fixes $\delta_{\text{CP}}$ to twelve significant
figures in the axiom; the measurement must converge toward $1.1296$\,rad
if UM is correct.

%% ─────────────────────────────────────────────────────────────────────────────
\section{The Full CKM Matrix}
%% ─────────────────────────────────────────────────────────────────────────────

\subsection{Wolfenstein parametrisation}

With $\lambda = 1/\varphi^3$ derived from the axiom and the
CP phase $\delta_{\text{CP}} = \arccos(W_B)$ fixed, the CKM matrix in
the standard Wolfenstein form~\cite{Wolfenstein1983} becomes

\begin{equation}
  V_{\text{CKM}}
  \approx
  \begin{pmatrix}
    1 - \lambda^2/2 & \lambda & A\lambda^3(\rho - i\eta) \\
    -\lambda & 1 - \lambda^2/2 & A\lambda^2 \\
    A\lambda^3(1-\rho-i\eta) & -A\lambda^2 & 1
  \end{pmatrix}
  + \mathcal{O}(\lambda^4),
  \label{eq:ckm}
\end{equation}
where $\lambda = 1/\varphi^3 = 0.23607$, $\delta_{\text{CP}} =
\arccos(W_B) = 1.1296$\,rad is encoded in the complex phases, and $A$,
$\rho$, $\eta$ are not yet derived from first principles within UM.

\subsection{The parameter $A$ and $|V_{cb}|$}

The entry $|V_{cb}| = A\lambda^2$ is governed by the $b \to c$
transition. In UM, this amplitude requires two simultaneous channel
contributions:

\begin{itemize}
  \item \textbf{Boundary channel} ($W_B$): the weak current crosses the
        boundary channel at the quark-$W$ vertex, mediating the flavour
        change.
  \item \textbf{Matter channel} ($W_M$): the two heavy-quark mass
        eigenstates couple through the matter channel, encoding the
        Higgs-sector insertion required when both quarks carry large
        masses.
\end{itemize}

The amplitude is the product of the two channel weights:
\begin{equation}
  |V_{cb}| = W_B \times W_M = 2r(1-r) \times r^2 = 2r^3(1-r).
  \label{eq:Vcb}
\end{equation}

\textbf{Numerical result.}
\begin{align}
  |V_{cb}|^{\text{UM}} &= 2\,(0.30902)^3\,(0.69098) \notag\\
  &= 2 \times 0.029524 \times 0.69098 = 0.040780.
\end{align}

PDG 2022: $|V_{cb}| = 0.04113$.
\[
  \text{Residual} = -0.85\%
  = 0.14\,n\,\varepsilon_{\text{floor}} \quad (n = 2).
\]

The two-channel product is an $n = 2$ quantity (two independent channel
weights, each contributing one crossing), so the accumulated floor is
$2\,\varepsilon_{\text{floor}} = 5.9\%$. The result is $0.14\,n\,\varepsilon$
below the floor.

The Wolfenstein parameter $A$ follows immediately:
\begin{equation}
  A = \frac{|V_{cb}|}{\lambda^2}
    = W_B \times W_M \times \varphi^6
    = 2r^3(1-r) \times \varphi^6.
  \label{eq:A}
\end{equation}

\begin{equation}
  A^{\text{UM}} = 0.040780 / 0.055728 = 0.7318.
\end{equation}

PDG 2022: $A = 0.7380$.
\[
  \text{Residual} = -0.85\%
  = 0.14\,n\,\varepsilon_{\text{floor}} \quad (n = 2).
\]

This completes the first-principles derivation of $A$ with no free
parameters.

\subsection{Comparison of CKM matrix elements}

Using $\lambda = 1/\varphi^3$, $A = W_B W_M \varphi^6 = 0.7318$, and the
PDG 2022 Wolfenstein parameters $\bar\rho = 0.122$, $\bar\eta = 0.355$:

\begin{center}
\begin{tabular}{lllll}
\toprule
Element & UM value & Observed (PDG 2022) & Residual & $n\,\varepsilon$ units \\
\midrule
$|V_{ud}|$ & $0.97213$ & $0.97373$ & $-0.16\%$ & $0.06$ \\
$|V_{us}|$ & $0.23607$ & $0.22450$ & $+5.15\%$ & $0.58\;(n{=}3)$ \\
$|V_{ub}|$ & $0.00370$ & $0.00361$ & $+2.49\%$ & $0.84$ \\
$|V_{cd}|$ & $0.23607$ & $0.22438$ & $+5.21\%$ & $0.59\;(n{=}3)$ \\
$|V_{cs}|$ & $0.97155$ & $0.97349$ & $-0.20\%$ & $0.07$ \\
$|V_{cb}|$ & $0.04078$ & $0.04113$ & $-0.85\%$ & $0.14\;(n{=}2)$ \\
$|V_{tb}|$ & $0.99914$ & $0.99912$ & $+0.00\%$ & $0.00$ \\
\bottomrule
\end{tabular}
\end{center}

All derived elements are within $1\,n\,\varepsilon_{\text{floor}}$ of
their observed values. The off-diagonal elements $|V_{us}|$ and $|V_{cd}|$
sit at $0.58$--$0.59\,n\,\varepsilon$ ($n = 3$), consistent with the QCD
correction anticipated above.

\subsection{The Jarlskog invariant}

The Jarlskog invariant $J$ is the unique rephasing-invariant measure of
CP violation~\cite{Jarlskog1985}. Using the UM-derived $A = 0.7318$ and
$\lambda = 1/\varphi^3$:
\begin{equation}
  J \approx A^2\,\lambda^6\,\bar\eta
  = (0.7318)^2 \times (0.23607)^6 \times 0.355
  \approx 3.29 \times 10^{-5}.
  \label{eq:J}
\end{equation}
The PDG 2022 value is $J = (3.08 \pm 0.15)\times 10^{-5}$; the residual
is $+6.8\%$. Note that $J$ uses the PDG values of $\bar\rho$ and $\bar\eta$
as inputs alongside the UM-derived $\lambda$ and $A$, so it is a
consistency check rather than a fully independent prediction.

%% ─────────────────────────────────────────────────────────────────────────────
\section{Discussion}
%% ─────────────────────────────────────────────────────────────────────────────

\subsection{Status of the two results}

The two results have very different precision:

\begin{itemize}
  \item $\delta_{\text{CP}} = \arccos(W_B) = 1.1296$\,rad sits at
        $0.31\,\varepsilon_{\text{floor}}$ from the PDG value. This is a
        sub-floor result requiring no further correction cycle. Within the
        geometric framework, the CP phase is exact.

  \item $\lambda = 1/\varphi^3 = 0.23607$ sits at $0.58\,n\,\varepsilon_{\text{floor}}$
        from the observed $|V_{us}|$ (accumulated floor $n = 3$). The formula
        $\lambda = 1/\varphi^3$ is the leading-order result; a sub-leading
        correction of order $r^2\lambda \sim 5\%$ is expected and consistent
        with the residual.
\end{itemize}

\subsection{Significance of the CP result}

The CKM CP-violating phase is one of the least constrained fundamental
constants. It has no theoretical derivation within the Standard Model; it
is simply measured. The UM derivation
\begin{equation}
  \delta_{\text{CP}} = \arccos\!\bigl(2r(1-r)\bigr)
  \label{eq:delta_clean}
\end{equation}
expresses it as a geometric angle in channel-weight space, depending only
on the single number $r = 1/(2\varphi)$. The 0.92\% agreement with the
PDG value is twelve times smaller than the current experimental
uncertainty. This makes $\delta_{\text{CP}}$ one of the sharpest tests
of the UM framework available to near-term experiments.

\subsection{Generation structure hierarchy}

The Wolfenstein hierarchy $\lambda \gg A\lambda^2 \gg A\lambda^3$ reflects
the three levels of inter-generation mixing: single-boundary, double, and
triple transitions respectively. In UM:
\begin{itemize}
  \item Single-boundary amplitude: $\lambda = 1/\varphi^3$.
  \item Double-boundary amplitude: $\lambda^2 = 1/\varphi^6$, suppressed
        by the probability of two sequential crossings.
  \item Triple-boundary amplitude: $\lambda^3 = 1/\varphi^9$, further
        suppressed, with a QCD factor $A$ encoding colour-sector corrections.
\end{itemize}

This cascade is geometrically natural in UM: each additional generation
step reduces the amplitude by $1/\varphi^3$, the same ratio that governs
single-generation mixing.

\subsection{What remains}

\begin{itemize}
  \item \textbf{Quark-sector QCD correction to $\lambda$.} Including the
        $\mathcal{O}(r^2)$ QCD correction should reduce the $\lambda$
        residual from $0.58$ to sub-floor. The dominant mechanism is the
        renormalisation of the boundary-crossing amplitude by the strong
        coupling.
  \item \textbf{PMNS matrix.} The neutrino mixing matrix may follow from
        an analogous boundary-crossing analysis in the lepton sector,
        with the lepton-sector depth (11 cycles vs.\ the quark-sector 4)
        producing the observed large mixing angles.
  \item \textbf{Precision test via Belle~II.} As $\delta_{\text{CP}}$
        measurements improve below $\pm 0.05$\,rad, the UM value
        $1.1296$\,rad will be directly testable.
  \item \textbf{$\bar\rho$ and $\bar\eta$ from first principles.} The
        Wolfenstein parameters $\bar\rho = 0.122$ and $\bar\eta = 0.355$
        are currently PDG inputs for the Jarlskog calculation; their
        derivation from the UM recursion (likely involving the $G_2$
        winding structure of Paper~6) will make $J$ a fully
        first-principles result.
\end{itemize}

%% ─────────────────────────────────────────────────────────────────────────────
\section{Prediction Table}
%% ─────────────────────────────────────────────────────────────────────────────

\begin{center}
{\color{gold}
\begin{tabular}{llllr}
\toprule
Observable & UM expression & UM value & Observed & Error\;($n\,\varepsilon$) \\
\midrule
$\lambda = |V_{us}|$
  & $1/\varphi^3$
  & $0.23607$ & $0.22450$ & $+5.15\%\;(0.58,\,n{=}3)$ \\
$\delta_{\text{CP}}$
  & $\arccos(W_B)$
  & $1.1296$\,rad & $1.14$\,rad & $-0.92\%\;(0.31,\,n{=}1)$ \\
$|V_{cb}|$
  & $W_B \times W_M$
  & $0.04078$ & $0.04113$ & $-0.85\%\;(0.14,\,n{=}2)$ \\
$A$
  & $W_B W_M \varphi^6$
  & $0.7318$ & $0.7380$ & $-0.85\%\;(0.14,\,n{=}2)$ \\
$J$ (mixed)
  & $A^2\lambda^6\bar\eta$
  & $3.29\times10^{-5}$ & $3.08\times10^{-5}$ & $+6.8\%$ \\
\bottomrule
\end{tabular}
}
\end{center}

%% ─────────────────────────────────────────────────────────────────────────────
\section*{Closing}
%% ─────────────────────────────────────────────────────────────────────────────

The CKM matrix parametrises the most complex structure in the electroweak
sector: three generations, four real parameters, one complex phase. The
Standard Model offers no explanation for any of them.

UM derives three of the four real parameters directly from the axiom: the
mixing angle $\lambda = 1/\varphi^3$, the transition amplitude
$|V_{cb}| = W_B \times W_M$, and the CP-violating phase
$\delta_{\text{CP}} = \arccos(W_B)$. All three sit within
$0.58\,n\,\varepsilon_{\text{floor}}$ of their observed values.

One axiom. Three channel weights. Three CKM parameters.

\begin{thebibliography}{9}
\bibitem{Wolfenstein1983}
L.~Wolfenstein,
\textit{Parametrization of the Kobayashi-Maskawa Matrix},
Phys.\ Rev.\ Lett.\ \textbf{51} (1983) 1945.

\bibitem{Jarlskog1985}
C.~Jarlskog,
\textit{Commutator of the Quark Mass Matrices in the Standard Electroweak Model and a Measure of Maximal CP Nonconservation},
Phys.\ Rev.\ Lett.\ \textbf{55} (1985) 1039.

\bibitem{PDG2022}
R.~L.~Workman \textit{et al.}\ (Particle Data Group),
\textit{Review of Particle Physics},
Prog.\ Theor.\ Exp.\ Phys.\ \textbf{2022}, 083C01 (2022).
\end{thebibliography}

\end{document}
